\subsubsection{Subsonic Inflow Boundary Condition}
\label{SEC:BC:Inflow:Subsonic}

We first need to introduce notation for some quantities that will be utilized
frequently when describing the different subsonic inflow methods.
%
For the velocity vector $\vec{u}$, let \vmag{} represent the total magnitude of
the vector, \thatis
%
\[ \vmag = \norm{\vec{u}} \]
%
We also introduce the following notation to represent the vector of the
directional cosines for each respective component of a velocity vector
%
\begin{equation}
  \dircos = \frac{\vec{u}}{\vmag}
  \label{EQ:BC:Inflow:Subsonic:DirCosines}
\end{equation}
%
where the accent ${\hat{\phantom{a}}}$ is used here to denote this is a unit
vector by definition.
%
For most of the following subsonic inflow methods, the directional cosines of
the prescribed inflow velocity are used to fix the angle of attack for the right
state velocity, therefore $\dircos_\subR=\dircos_\subBC$.

%%%%%%%%%%%%%%%%%%%%%%%%%%%%%%%%%%%%%%%%%%%%%%%%%%%%%%%%%%%%%%%%%%%%%%%%%%%%%%%%

Using these definitions, the velocity in the direction of the
(outward pointing) surface normal can be written as
%
\begin{equation}
  v_n = \vec{u}\cdot\normal
      = \left(\vmag\dircos\right)\cdot\normal
      = \vmag\left(\dircos\cdot\normal\right)
      = \vmag\dcnorm
  \label{EQ:BC:Inflow:Subsonic:NormalVelocity}
\end{equation}
%
where
\vmag{} is the magnitude of the velocity vector defined above and
\dcnorm{} represents the fraction of \vmag{} in the direction of \normal.
%
For a flow that is perpendicular to the surface, $\dcnorm{} = -1$.
%
We also note that $v_n < 0$ on an inflow surface and
$\abs{v_n} < a$ on a subsonic inflow.


%%%%%%%%%%%%%%%%%%%%%%%%%%%%%%%%%%%%%%%%%%%%%%%%%%%%%%%%%%%%%%%%%%%%%%%%%%%%%%%%

The Riemann invariants and corresponding characteristic velocities are
given by
%
\begin{align}
      R^+       &= v_n + \frac{2a_\subL}{\gamma-1}
  &   \lambda_1 &= v_n + a > 0
  &&  \text{Outgoing}
  \\
      R^-       &= v_n - \frac{2a_\subL}{\gamma-1}
  &   \lambda_2 &= v_n - a < 0
  &&  \text{Incoming}
\end{align}

%%%%%%%%%%%%%%%%%%%%%%%%%%%%%%%%%%%%%%%%%%%%%%%%%%%%%%%%%%%%%%%%%%%%%%%%%%%%%%%%

It should be noted that for multi-component gas mixtures, the right state is
determined in a manner similar to method 4 (\ref{SEC:BC:Inflow:Subsonic:4})
unless method 2 (\ref{SEC:BC:Inflow:Subsonic:2}) is specified via
\inp[nmlBC::SubInflowMethod]{Sub\_Inflow\_Method}.

\devNotes{Need to fix subsonic inflow angle-of-attack bug for TPG subroutine.}

%%%%%%%%%%%%%%%%%%%%%%%%%%%%%%%%%%%%%%%%%%%%%%%%%%%%%%%%%%%%%%%%%%%%%%%%%%%%%%%%

\subsubsubsection{\texorpdfstring{\inptext{Sub\_Inflow\_Method} = 1}%
                                 {Sub\_Inflow\_Method = 1}}
\label{SEC:BC:Inflow:Subsonic:1}

This method performs a Newton iteration on the total temperature to get the
right state.
%
It is based on the method found on pages 8--9 of \linecite{Cooper_1990} from
the subsection titled \emph{Internal Flows} in Section 4.2.1.
%
This method uses the specified boundary condition values $p_{t,\subBC}$,
$T_{t,\subBC}$, and $\dircos_\subBC$ along with the outward propagating
Riemann invariant, $R^+$, to determine the right state.

%%%%%%%%%%%%%%%%%%%%%%%%%%%%%%%%%%%%%%%%%%%%%%%%%%%%%%%%%%%%%%%%%%%%%%%%%%%%%%%%

We start with enforcing that the total temperature for the right state is equal
to the prescribed total temperature
%
\[
  T_{t,\subR} = T_\subR \left(1 + \frac{\gamma-1}{2}M^2_\subR\right)
              = T_{t,\subBC}
\]
%
Using the relations $M=\vmag/a$, $a^2=\gamma\rgas T$, and
$c_p=\gamma\rgas/\left(\gamma-1\right)$, the total temperature can be rewritten
as
%
\begin{equation}
  T_{t,\subBC} = T_\subR + \frac{1}{2c_p}\vmag^2_\subR
  \label{EQ:BC:Inflow:Subsonic1:Ttotal}
\end{equation}
%
The next step is to set the outgoing Riemann invariant, $R^+$, using the
interior state.
%
We use the subscript $\subL$ to note that the invariant is computed from
the left (interior) state.
%
\begin{equation}
  R^+_\subL = \vec{u}_\subL\cdot\normal + \frac{2a_\subL}{\gamma-1}
  \label{EQ:BC:Inflow:Subsonic1:RplusL}
\end{equation}
%
We want to find a right state normal velocity and speed of sound that keeps the
outgoing invariant constant, or
%
\begin{equation}
  \vec{u}_\subR\cdot\normal + \frac{2a_\subR}{\gamma-1} = R^+_\subL
  \label{EQ:BC:Inflow:Subsonic1:Rplus}
\end{equation}
%
Using \eref{EQ:BC:Inflow:Subsonic:NormalVelocity} to represent the velocity
normal to the surface, we can rearrange \eref{EQ:BC:Inflow:Subsonic1:Rplus} as
%
\begin{equation}
  \vmag_\subR = \frac{1}{\dcnorm_\subBC}
                \left( R^+_\subL - \frac{2a_\subR}{\gamma-1} \right)
  \label{EQ:BC:Inflow:Subsonic1:vmagR}
\end{equation}
%
where the subscript \subBC{} has been added to $\dcnorm_\subBC$ to indicate
that we are using the direction cosines from the boundary condition velocity
vector to enforce the right state velocity to the prescribed angles of attack.
%
Finally, we can combine \eref{EQ:BC:Inflow:Subsonic1:Ttotal} and
\eref{EQ:BC:Inflow:Subsonic1:vmagR} resulting in the equation
%
\begin{equation*}
  T_{t,\subBC} = T_\subR + \frac{1}{2c_p}
    \left[ 
      \frac{1}{\dcnorm_\subBC}
      \left( R^+_\subL - \frac{2a_\subR}{\gamma-1} \right)
    \right]^2
\end{equation*}
%
With $a_\subR=\sqrt{\gamma\rgas T_\subR}$, this equation becomes solely
dependent on the right state temperature $(T_\subR)$ with all the other
variables being either constant or prescribed by the boundary condition state
%
\begin{equation}
  f\left(T_\subR\right) = T_\subR + \frac{1}{2c_p}
    \left[ 
      \frac{1}{\dcnorm_\subBC}
      \left( R^+_\subL - \frac{2a_\subR}{\gamma-1} \right)
    \right]^2 - T_{t,\subBC}
  = 0
  \label{EQ:BC:Inflow:Subsonic1:fTr}
\end{equation}
%
We can solve this numerically using a Newton iteration
%
\begin{align*}
  T^{n+1}_\subR =&
         T^n_\subR - \frac{f\left(T^n_\subR\right)}{f'\left(T^n_\subR\right)}
  \\
  a^{n+1}_\subR =& \sqrt{ \gamma \rgas T^{n+1}_\subR }
\end{align*}
%
where $n$ is the current iteration of the Newton solver.
%
All that remains is taking the derivative of $f\left(T_\subR\right)$, which
condenses nicely to
%
\begin{equation*}
  f'\left(T_\subR\right) = 1 - \frac{\vmag_\subR}{\dcnorm_\subBC a_\subR}
\end{equation*}
%
After the Newton iteration converges to a solution for $T_\subR$ and
$a_\subR$,
%the speed of sound is updated using $a_\subR=\sqrt{\gamma\rgas T_\subR}$,
we can then continue on to compute the final right state variables needed
for the boundary condition.
%
The next step is to update $\vmag_\subR$ using
\eref{EQ:BC:Inflow:Subsonic1:vmagR} from which a Mach number,
$M_\subR=\vmag_\subR / a_\subR$, can be evaluated.
%
Finally, the primitive variables for the right state can be computed as
%
\begin{alignat}{1}
  \vec{u}_\subR & = \vmag_\subR \dircos_\subBC
  \label{EQ:BC:Inflow:Subsonic1:velR}
  \\
        p_\subR & = p_{t,\subBC}
                    \left( 1 + \frac{\gamma-1}{2} M^2_\subR
                    \right)^{-\gamma / \left( \gamma - 1 \right)}
  \label{EQ:BC:Inflow:Subsonic1:presR}
  \\
     \rho_\subR & = \frac{p_\subR}{\rgas T_\subR}
  \label{EQ:BC:Inflow:Subsonic1:densR}
\end{alignat}

%%%%%%%%%%%%%%%%%%%%%%%%%%%%%%%%%%%%%%%%%%%%%%%%%%%%%%%%%%%%%%%%%%%%%%%%%%%%%%%%

\subsubsubsection{\texorpdfstring{\inptext{Sub\_Inflow\_Method} = 2}%
                                 {Sub\_Inflow\_Method = 2}}
\label{SEC:BC:Inflow:Subsonic:2}

This method is exactly the opposite of the subsonic outflow method, using the
density and velocity specified for the boundary condition and taking the static
pressure from the interior
%
\begin{alignat*}{1}
  \rho_\subR & = \rho_\subBC \\
  \vec{u}_\subR & = \vec{u}_\subBC \\
  p_\subR & = p_\subL
\end{alignat*}
%
This method does not enforce the total temperature or pressure at the inflow, so
caution should be taken with its use.
%
This method is probably only suitable for a subsonic inflow that is far upstream
from the body of interest. Even then, it is probably best to use one of the
other subsonic inflow methods or the characteristic boundary condition
(\ref{SEC:BC:Farfield:Characteristic}) instead.
%

%%%%%%%%%%%%%%%%%%%%%%%%%%%%%%%%%%%%%%%%%%%%%%%%%%%%%%%%%%%%%%%%%%%%%%%%%%%%%%%%

\subsubsubsection{\texorpdfstring{\inptext{Sub\_Inflow\_Method} = 3}%
                                 {Sub\_Inflow\_Method = 3}}
\label{SEC:BC:Inflow:Subsonic:3}

This method uses the specified boundary condition values $p_{t,\subBC}$,
$\dircos_\subBC$, and $H_{t,\subBC}$ (total enthalpy) along with the outward
propagating Riemann invariant, $R^+$, to determine the right state.
%
It is very similar to
$\inptext{Sub\_Inflow\_Method}=1$ (\ref{SEC:BC:Inflow:Subsonic:1}),
but a quadratic equation is solved for $a_\subR$
instead of using a Newton iteration for $T_\subR$.
%
It is also equivalent to the method found in Section 2.7 of
\linecite{Carlson_2011} (pages 7--9), though somewhat different notation is
used there and the flow is assumed to be normal to the boundary.

We start by rewriting \eref{EQ:BC:Inflow:Subsonic1:fTr} in terms of the
speed of sound, with $R^+$ computed from the left (interior) state using
\eref{EQ:BC:Inflow:Subsonic1:RplusL},
\begin{equation}
  \frac{a_\subR^2}{\gamma \rgas} + \frac{1}{2c_p}
    \left[ 
      \frac{1}{\dcnorm_\subBC}
      \left( R^+_\subL - \frac{2a_\subR}{\gamma-1} \right)
    \right]^2 - T_{t,\subBC}
  = 0
\end{equation}
%
and then rearranging the terms into a standard quadratic form
with $H_{t,\subBC} = c_p T_{t,\subBC}$.
\begin{equation}
    \left[ \dcnorm^2_\subBC + \frac{2}{\gamma-1} \right] a^2_\subR
  - \left[ 2 R^+_\subL \right] a_\subR
  + \frac{\gamma-1}{2}
    \left[ \left( R^+_\subL \right)^2 - 2 \dcnorm^2_\subBC H_{t,\subBC} \right]
  = 0
% \label{EQ:BC:Inflow:Subsonic3:Quadratic}
\end{equation}
%
Using the quadratic formula to solve this equation results in two solutions for
$a_\subR$.
%
If for some reason both solutions are negative, the speed of sound from the
interior state is used for the right state, \thatis $a_\subR = a_\subL$.
%
Otherwise $a_\subR$ is taken to be the maximum value between the two solutions.

%%%%%%%%%%%%%%%%%%%%%%%%%%%%%%%%%%%%%%%%%%%%%%%%%%%%%%%%%%%%%%%%%%%%%%%%%%%%%%%%

The temperature can be computed from $T_\subR = a_\subR^2 / \gamma \rgas$,
and the remaining primitive variables can be updated in the same way as in
\inptext{Sub\_Inflow\_Method}=1 (\sref{SEC:BC:Inflow:Subsonic:1}).
%
$\vmag_\subR$ is updated using \eref{EQ:BC:Inflow:Subsonic1:vmagR}
from which a Mach number, $M_\subR=\vmag_\subR / a_\subR$, can be evaluated.
%
Then $\vec{u}_\subR$, $p_\subR$, and $\rho_\subR$ are computed from
\eref{%
EQ:BC:Inflow:Subsonic1:velR,%
EQ:BC:Inflow:Subsonic1:presR,%
EQ:BC:Inflow:Subsonic1:densR%
}.

%%%%%%%%%%%%%%%%%%%%%%%%%%%%%%%%%%%%%%%%%%%%%%%%%%%%%%%%%%%%%%%%%%%%%%%%%%%%%%%%

\subsubsubsection{\texorpdfstring{\inptext{Sub\_Inflow\_Method} = 4}%
                                 {Sub\_Inflow\_Method = 4}}
\label{SEC:BC:Inflow:Subsonic:4}

This is the default method used for the subsonic inflow boundary condition,
because it is the simplest of the four available methods that also enforces the
total conditions specified for the boundary condition.
%
This method uses the specified boundary condition values $p_{t,\subBC}$,
$T_{t,\subBC}$, and $\dircos_\subBC$ along with the interior Mach number at the
boundary to determine the right state.
%
Since both $p_{t,\subBC}$ and $T_{t,\subBC}$ are specified, this implies that
$\rho_{t,\subBC}$ is also specified through the ideal gas law.

%%%%%%%%%%%%%%%%%%%%%%%%%%%%%%%%%%%%%%%%%%%%%%%%%%%%%%%%%%%%%%%%%%%%%%%%%%%%%%%%

It starts by enforcing that the right state Mach number is set to the interior
Mach number
\[
  M_\subR = M_\subL = \frac{\vmag_\subL}{a_\subL}
\]
%
which then allows us to compute the static temperature for the
right state based on the specified total temperature using
%
\[
  T_\subR = T_{t,\subBC} \left(1 + \frac{\gamma-1}{2}M^2_\subR\right)^{-1}
\]
%
From here, the magnitude of the right state velocity can be computed as
%
\[
  \vmag_\subR = M_\subR a_\subR = M_\subR \sqrt{\gamma \rgas T_\subR}
\]
%
which gives us everything that is needed to finally compute the primitive
variables for the right state
%
\begin{alignat}{1}
     \rho_\subR & = \left( \frac{p_{t,\subBC}}{\rgas T_{t,\subBC}} \right)
                    \left( 1 + \frac{\gamma-1}{2} M^2_\subR
                    \right)^{-1 / \left( \gamma - 1 \right)} \\
  \vec{u}_\subR & = \vmag_\subR \dircos_\subBC \\
        p_\subR & = \rho_\subR \rgas T_\subR
  \label{EQ:BC:Inflow:Subsonic4:PV2}
\end{alignat}

%%%%%%%%%%%%%%%%%%%%%%%%%%%%%%%%%%%%%%%%%%%%%%%%%%%%%%%%%%%%%%%%%%%%%%%%%%%%%%%%

%If looking at this from a practical stance separate of the GFR code, most users
%would want to specify a total pressure and total temperature for a subsonic
%inflow boundary condition.
%%
%In this case, you could think of this as first computing $p_\subR$ from
%\eref{EQ:BC:Inflow:Subsonic3:PV} using the boundary condition total pressure,
%$p_{t,\subBC}$, and then using the ideal gas law to compute $\rho_\subR$ from
%$p_\subR$ and $T_\subR$.
%%
%Either way, the right state variables in \eref{EQ:BC:Inflow:Subsonic4:PV}
%should come out the same assuming $\rho_{t,\subBC}$, $p_{t,\subBC}$, and
%$T_{t,\subBC}$ all conform to the ideal gas law.
%%
%Within GFR, this does not really matter since the requirement from
%\sref{SEC:BC:Specify:FullState} of having to fully define each flow-through
%boundary condition means these three total conditions are all known and
%consistent.

%%%%%%%%%%%%%%%%%%%%%%%%%%%%%%%%%%%%%%%%%%%%%%%%%%%%%%%%%%%%%%%%%%%%%%%%%%%%%%%%

%###############################################################################

